% Each hemisphere sees the opposite half of the visual field. Build: sh tex/tex2svg.sh (from docs/)
\documentclass[tikz, border=0pt]{standalone}
\usepackage{neuroai-fig}
% local red/blue pair: left vs right half are easier to tell apart than teal/coral
\definecolor{redF}{HTML}{FCEBEB}  \definecolor{redS}{HTML}{A32D2D}  \definecolor{redD}{HTML}{791F1F}  \definecolor{arrR}{HTML}{E24B4A}
\definecolor{blueF}{HTML}{E6F1FB} \definecolor{blueS}{HTML}{185FA5} \definecolor{blueD}{HTML}{0C447C} \definecolor{arrB}{HTML}{378ADD}
\begin{document}
% coordinates: 1 unit = 0.75pt, y grows downwards (seen from above, scene at the top)
\begin{tikzpicture}[figure, x=0.75pt, y=-0.75pt, every node/.append style={align=center}]
\tikzset{
  ray/.style={draw=#1, line width=0.75pt, opacity=0.7},
  fiber/.style={draw=#1, line width=1.9pt},
  retina/.style={draw=#1, line width=4.5pt},
  field/.style={box=#1, align=center, minimum width=150pt, minimum height=37pt},
}
% fixation line and point
\draw[draw=lobe, line width=0.75pt, dash pattern=on 2.2pt off 3pt] (340,40) -- (340,600);
\node[text=note, above=2pt] at (340,40) {You look here};
\fill[arrG] (340,65) circle (3.7pt);
% the two halves of the scene
\node[field=red, text=redD] (left)  at (200,65) {\large Left half of scene};
\node[field=blue, text=blueD] (right) at (480,65) {\large Right half of scene};

% eyes: both turn to face the fixation point, so it lands on each fovea (the split between
% the two retina halves). A scene point enters the pupil and lands on the opposite side of
% the fovea. Pupil on the eyeball surface: a ray with direction t lands at angle 2t+180-g.
\coordinate (F) at (340,65);
\foreach \ex/\name in {240/leye, 440/reye} {
  \path let \p1 = ($(F)-(\ex,200)$), \n1 = {atan2(\y1,\x1)} in \pgfextra{\xdef\g{\n1}};  % gaze angle
  \coordinate (P) at ($(\ex,200)+(\g:30pt)$);                                             % pupil
  \path let \p1 = ($(P)-(left.east)$), \p2 = ($(P)-(right.west)$),
            \n1 = {atan2(\y1,\x1)}, \n2 = {atan2(\y2,\x2)} in \pgfextra{\xdef\tl{\n1}\xdef\tr{\n2}};
  \draw[draw=lobe, line width=0.75pt, dash pattern=on 2.2pt off 3pt] (F) -- ($(\ex,200)+({\g+180}:30pt)$);
  \draw[ray=arrR] (left.east) -- (P) -- ($(\ex,200)+({2*\tl+180-\g}:30pt)$);
  \draw[ray=arrB] (right.west) -- (P) -- ($(\ex,200)+({2*\tr+180-\g}:30pt)$);
  \draw[draw=note, line width=0.75pt] (\ex,200) circle (30pt);
  \draw[retina=arrB] ($(\ex,200)+({\g+90}:30pt)$) arc[start angle=\g+90, end angle=\g+180, radius=30pt];
  \draw[retina=arrR] ($(\ex,200)+({\g+180}:30pt)$) arc[start angle=\g+180, end angle=\g+270, radius=30pt];
  \fill[note] (P) circle (3pt);
  \coordinate (\name-T) at ($(\ex,200)+({\g+135}:30pt)$);  % fibre starts: middle of each retina half
  \coordinate (\name-C) at ($(\ex,200)+({\g+225}:30pt)$);
}
\node[anchor=east, font=\large\bfseries, text=ink] at (190,200) {Left eye};
\node[anchor=west, font=\large\bfseries, text=ink] at (490,200) {Right eye};

% hemispheres
\node[field=blue, minimum width=0pt] (llgn) at (210,431) {\textcolor{blueD}{\large Left LGN（外侧膝状体）}\\ Relays to V1};
\node[field=red, minimum width=0pt] (rlgn) at (470,431) {\textcolor{redD}{\large Right LGN（外侧膝状体）}\\ Relays to V1};
\node[field=blue, minimum width=0pt] (lv1) at (210,523) {\textcolor{blueD}{\large Left V1（初级视觉皮层）}\\ Gets right half only};
\node[field=red, minimum width=0pt] (rv1) at (470,523) {\textcolor{redD}{\large Right V1（初级视觉皮层）}\\ Gets left half only};
\draw[arr=arrB] (llgn) -- (lv1);
\draw[arr=arrR] (rlgn) -- (rv1);
\begin{pgfonlayer}{lobes}
  \node[lobe, inner sep=12pt, fit={(llgn) (lv1) ($(lv1.south)+(0,22)$)}] (lhem) {};
  \node[lobe, inner sep=12pt, fit={(rlgn) (rv1) ($(rv1.south)+(0,22)$)}] (rhem) {};
\end{pgfonlayer}
\node[anchor=south west, text=note] at ($(lhem.south west)+(8pt,5pt)$) {Left hemisphere};
\node[anchor=south west, text=note] at ($(rhem.south west)+(8pt,5pt)$) {Right hemisphere};

% optic nerve fibres: the nasal half of each retina crosses at the chiasm
\draw[fiber=arrB] (leye-T) .. controls (212,300) and (200,340) .. ([xshift=-12pt]llgn.north);
\draw[fiber=arrR] (leye-C) .. controls (340,320) .. ([xshift=-12pt]rlgn.north);
\draw[fiber=arrB] (reye-T) .. controls (340,320) .. ([xshift=12pt]llgn.north);
\draw[fiber=arrR] (reye-C) .. controls (468,300) and (480,340) .. ([xshift=12pt]rlgn.north);
% label sits in the free wedge right of the crossing, left of the right eye's own fibre
\node[anchor=west, inner sep=0pt, align=left, font=\small\bfseries, text=ink] at (380,318) {Optic chiasm\\[1pt]（视交叉）};

\node[text=note] at (340,618) {Seen from above, both eyes turned toward the fixation point};
\end{tikzpicture}
\end{document}
